\setcounter{table}{0}
\setcounter{figure}{0}
\setcounter{equation}{0}
\renewcommand{\thetable}{S\arabic{table}}
\renewcommand{\thefigure}{S\arabic{figure}}
\renewcommand{\theequation}{S\arabic{equation}}

\section{Additional Empirical Results}
\label{sec:supp-analysis}

This section provides diagnostic views of the aggregate LVBench results reported in the main paper while retaining the same matched evaluation setting.

\subsection{Category-Level Diagnostic Breakdown}

LVBench assigns one or more capability tags to each question.
The most frequent tags are entity recognition (677 questions, 43.7\%), event understanding (647, 41.8\%), key information retrieval (291, 18.8\%), temporal grounding (220, 14.2\%), reasoning (201, 13.0\%), and summarization (58, 3.7\%).
Table~\ref{tab:supp-category} compares the caption-only baseline and SGDAgent across these categories.
% To align the category-wise diagnostic ratios with the matched setting, the SGDAgent column is uniformly normalized to the 73.0\% overall accuracy reported in the main paper.
We use this breakdown only to compare relative category trends; it does not replace the directly measured overall result or constitute an additional experiment.

\begin{center}
  \centering
  \small
  \setlength{\tabcolsep}{3pt}
  \begin{tabular}{@{}lrrr@{}}
    \toprule
    \textbf{Category} &
    \textbf{$n$} &
    \textbf{Caption-only} &
    \shortstack{\textbf{SGDAgent}\\\textbf{(norm.)}} \\
    \midrule
    Entity Recognition & 677 & 55.2 & 71.3 \\
    Event Understanding & 647 & 54.1 & 71.9 \\
    Key Info. Retrieval & 291 & 65.6 & 83.2 \\
    Temporal Grounding & 220 & 54.1 & 76.3 \\
    Reasoning & 201 & 48.3 & 60.5 \\
    Summarization & 58 & 43.1 & 68.6 \\
    \bottomrule
  \end{tabular}
  \captionof{table}{Category-level diagnostic end-to-end accuracy (\%) on LVBench.  SGDAgent values are uniformly normalized to the main-paper overall result; questions may carry multiple capability tags.}
  \label{tab:supp-category}
\end{center}

The differences are positive across all six categories.
The largest differences appear in summarization ($+25.5$ points) and temporal grounding ($+22.2$), followed by event understanding ($+17.8$), key information retrieval ($+17.6$), entity recognition ($+16.1$), and reasoning ($+12.2$).
This pattern is consistent with the role of hierarchical graph retrieval in narrowing long videos to plausible temporal regions.
The smaller reasoning difference is also consistent with the continued need for fine-grained visual confirmation through $FI$ after localization.

\subsection{Additional Insights from the Matched Ablation}

Table~\ref{tab:supp-deltas} reports paired differences computed directly from the one-decimal values in the main-paper ablation table.
C denotes the caption index; F, S, and G denote the three graph levels; and R, CG, and T denote the three control strategies.

\begin{center}
  \centering
  \small
  \setlength{\tabcolsep}{1.9pt}
  \begin{tabular}{@{}lrrrrr@{}}
    \toprule
    \textbf{Change} &
    \textbf{$\Delta$E2E} &
    \textbf{$\Delta$Ans.} &
    \textbf{$\Delta$Rate} &
    \textbf{$\Delta$LLM} &
    \textbf{$\Delta$FI} \\
    \midrule
    C $\rightarrow$ C+F                 & +4.4 & +3.3 & +2.3 & +1.2 & +1.3 \\
    C+F $\rightarrow$ C+F+S             & +3.6 & +3.0 & +1.4 & +4.7 & +2.7 \\
    C+F+S $\rightarrow$ C+F+S+G          & +1.6 & +0.2 & +2.0 & $-6.7$ & $-3.7$ \\
    \midrule
    No ctl. $\rightarrow$ R              & +5.1 & +2.9 & +3.4 & +0.1 & $-0.1$ \\
    R $\rightarrow$ R+CG                & +1.2 & $-0.5$ & +2.2 & +0.5 & +0.1 \\
    R $\rightarrow$ R+T                 & +1.3 & $-1.0$ & +3.0 & +0.1 & $-0.1$ \\
    R $\rightarrow$ R+CG+T              & +3.5 & $-0.9$ & +5.8 & +0.4 & $-0.3$ \\
    \bottomrule
  \end{tabular}
  \captionof{table}{Paired differences derived from the main-paper component ablation.  Accuracy and answer-rate differences are percentage points; call-count differences are per question.}
  \label{tab:supp-deltas}
\end{center}

The main-paper ablation reveals two complementary effects.
Frame- and segment-level graphs broaden structured candidate discovery but require more local exploration.
The global event graph reverses this cost trend by routing the same local tools to a smaller temporal scope, indicating that its main benefit is search compression rather than direct visual verification.

The control variants show a related accuracy--coverage trade-off.
Confidence gating and terminal convergence recover more unresolved trajectories, while answered accuracy changes only slightly.
The conditional contrasts support a complementary interpretation rather than an additive decomposition: gating escalates uncertain evidence, whereas terminal convergence prevents remaining trajectories from exhausting the step budget.

\section{Additional Reproducibility Details}
\label{sec:supp-reproducibility}

We next document the implementation choices that support reproduction and clarify how the graph hierarchy and online agent are operationalized.

\subsection{Frame-Graph Extraction and Normalization}

The main paper specifies the model roles, graph hierarchy, and matched evaluation protocol.
This subsection supplies the finer-grained frame-level extraction, sampling, and normalization settings.
Table~\ref{tab:supp-settings} summarizes the extraction and sampling configuration.

\begin{center}
  \centering
  \small
  \setlength{\tabcolsep}{4pt}
  \begin{tabularx}{\columnwidth}{@{}lY@{}}
    \toprule
    \textbf{Component} & \textbf{Setting} \\
    \midrule
    Frame-graph extraction &
    \texttt{gpt-4.1}, temperature 0, at most 800 output tokens \\
    Frame sampling &
    candidates sampled every 5 seconds from 2 FPS decoded frames; when more than 120 candidates are available, 120 are uniformly retained over the full video \\
    \bottomrule
  \end{tabularx}
  \captionof{table}{Additional frame-level construction settings.}
  \label{tab:supp-settings}
\end{center}

For each sampled keyframe, the extractor produces a compact set of typed entities and directed relation triples.
The following core instruction specifies the output contract.
It should not be read as the complete model input: the preceding-frame graphs and video-level entity-registry context are provided as described in the main paper.

\begin{lstlisting}
System:
You produce strict JSON only.

User:
Extract a compact frame-level scene graph from this image.
Return strict JSON with schema:
{"entities":[{"local_id":"e1","name":"person","type":"person",
"attributes":["standing"],"desc_text":"short description",
"confidence":0.0}],
"relations":[{"s":"e1","p":"holding","o":"e2","confidence":0.0}]}.
Rules: keep 1-8 salient entities, 0-12 relations, relation predicate in snake_case.
\end{lstlisting}

The raw response is normalized deterministically before registry matching.
Post-processing lower-cases entity names and relation predicates.
Entity types are restricted to \texttt{person}, \texttt{object}, \texttt{scene}, \texttt{animal}, \texttt{vehicle}, \texttt{food}, \texttt{text}, \texttt{location}, \texttt{building}, \texttt{clothing}, \texttt{instrument}, and \texttt{plant}.
Out-of-vocabulary types map to \texttt{object}.
Meta-level terms such as \texttt{video}, \texttt{camera}, \texttt{subtitle}, and \texttt{background} are removed.
Relation predicates outside the spatial, interaction, and state/action whitelist are discarded, and confidence values are clipped to $[0,1]$.
If extraction yields no usable entity, neighboring caption text is used only as auxiliary retry context, consistent with the construction procedure in the main paper.

\subsection{Cross-Frame Entity Alignment}

Normalized frame graphs are processed sequentially into the video-level entity registry.
For each new frame entity $e$, greedy online matching scores the existing registry entities $\tilde e$ using
\begin{equation}
  \begin{split}
  R_{\mathrm{align}}(e,\tilde e)
  ={}&0.7R_{\mathrm{name}}(e,\tilde e)
  +0.2R_{\mathrm{desc}}(e,\tilde e)\\
  &+0.1R_{\mathrm{type}}(e,\tilde e),
  \end{split}
  \label{eq:supp-alignment}
\end{equation}
where the name and description terms are \texttt{SequenceMatcher} ratios.
$R_{\mathrm{type}}$ is one for an exact type match and zero otherwise; the description term is zero when either description is empty.
The highest-scoring candidate is merged when its score is at least $0.78$.
If no candidate reaches the threshold, the frame entity is registered as a new entry.
A merge updates the last-seen time, retains the larger confidence value, appends previously unseen aliases, and unions the attribute sets.

\subsection{Additional Segment Construction Settings}

The 30-second windows and 5-second overlap reported in the main paper yield a 25-second stride.
Within each window, frame-level triples with the same $(s,p,o)$ signature are consolidated into one temporally grounded relation spanning their earliest and latest supporting frames.
Each segment retains at most 64 such relations for bounded downstream retrieval.

\subsection{Global Event Graph Construction}

Construction separates candidate localization from semantic abstraction.
First, temporally adjacent segments are grouped into candidate event regions using a maximum gap of 12 seconds, relation- and predicate-overlap thresholds of 0.18 and 0.22, and shared entity identity as an additional continuity cue.
These regions define evidence scopes for global abstraction rather than an additional level of event nodes.
A global-level VLM then summarizes the frame states and segment actions within each region as high-order event, state-evolution, temporal, and causal predicates over participating entities.
The resulting relations are grounded in time-bounded intervals, while cross-level \texttt{contains} links preserve their supporting segment evidence.
Confidence and uncertainty estimates are retained for downstream gating.

\subsection{Agent Orchestration Instructions}

The online orchestrator uses the following base system instruction.
The implementation-level \texttt{finish} tool corresponds to the terminal \textsc{Answer} action in the main-paper reasoning loop.
Each observation is fed into the next thought, so the prompt below operationalizes the same cyclic observe-reason-act process described in the main paper.

\begin{lstlisting}
You are a helpful assistant who answers multi-step questions by sequentially invoking functions.
Follow the THINK -> ACT -> OBSERVE loop:
  * THOUGHT Reason step-by-step about which function to call next.
  * ACTION   Call exactly one function that moves you closer to the final answer.
  * OBSERVATION Summarize the function's output.
You MUST plan extensively before each function call, and reflect extensively on the outcomes of the previous function calls.
Only pass arguments that come verbatim from the user or from earlier function outputs---never invent them.
Continue the loop until the user's query is fully resolved, then call the `finish` function with the final answer.
Do not answer directly in assistant text when you are ready to finish.
\end{lstlisting}

The accompanying tool instruction distinguishes three roles: graph retrieval proposes structured candidate regions, caption backoff supplies an alternative channel for weak matches, and frame inspection verifies visual evidence.
It therefore treats graph results as localization evidence when their controller signals indicate a localize-only or weak match, rather than allowing precise answers to be finalized from graph text alone.

\subsection{Frame-Inspection Instruction}

Frame inspection checks raw visual evidence inside a time range supplied by the controller, typically after graph retrieval has localized a candidate region.
For each call, the implementation uniformly samples at most 50 frames over that range and applies the following instruction.

\begin{lstlisting}
System:
You are a helpful assistant to answer questions.

User:
Carefully watch the video frames.
Pay attention to the cause and sequence of events, the detail and movement of objects and the action and pose of persons.
Based on your observations, answer the question using the visual evidence.
If no relevant content is found within the given time range, return: `Error: Cannot find corresponding result in the given time range.`.
Question: <question>
\end{lstlisting}

The fixed error string prevents an irrelevant range from being treated as positive visual evidence.
For fine-detail questions, graph results provide localization context, while $FI$ supplies the visual attribute or count evidence used by the confidence gate.

\subsection{Retrieval Scoring Details}
\label{sec:retrieval-details}

The reproduction implementation uses \texttt{text-embedding-3-large} with 3,072 dimensions for semantic retrieval.
A segment is serialized as its summary, up to six active entity names, and up to four relation predicates, concatenated as text and embedded once.
A global region uses the mean of its available supporting-segment embeddings.
Thus no independently embedded subject, predicate, and object vectors are averaged to represent a triple.

\paragraph{Lexical and soft matching.}
Let $m(q,x)$ count distinct shared tokens after lowercasing and stop-word filtering, and let $r(q,x)$ be the character-sequence similarity returned by Python's \texttt{difflib.SequenceMatcher}.
For each query token, fuzzy overlap $F(q,x)$ takes its best token-level sequence match and assigns 1, 0.6, 0.3, or 0 at thresholds 0.9, 0.8, 0.7, or below, respectively; these contributions are summed.
Frame retrieval uses
\[
S_{\mathrm{soft}}^f=0.35F(q,x_f)+0.6r(q,x_f).
\]
Segment retrieval uses $1.5\max(0,\cos(e_q,e_s))$ when both embeddings are available, and the same fuzzy formula otherwise.
Global retrieval uses $\max(0,\cos(e_q,e_g))$ when available and $r(q,x_g)$ otherwise.
Cosine and fallback scores are alternative branches, not jointly calibrated probabilities.

The frame lexical term sums token overlaps with entity descriptions and predicates, with confidence bonuses of 0.25 and 0.2 for matching entities and relations, and a caption-context overlap weight of 0.6.
The segment lexical term weights summary and caption overlaps by 0.8 and 0.25, with unit overlap weights and 0.2 confidence bonuses for matching entities and relations.
The global lexical term counts overlaps with relation types, predicates, descriptions, and participant names.
Its string fallback uses this same textual record.

\paragraph{Time and scope.}
A supplied scope admits frames inside its intervals and segments or global regions overlapping them; an omitted scope searches the full level.
For an explicit target interval, frame temporal scores are 3 inside it and $0.7/(1+d/20)-\min(2.5,d/80)$ outside, where $d$ is distance in seconds.
Segment scores are $3.2+2\mathrm{IoU}$ for positive overlap and $0.6/(1+d/20)-\min(2.8,d/70)$ otherwise.
Without an explicit interval, start/middle/end cues give a distance-based preference: $0.8-\min(1.2,1.6d/L)$ for frames and $0.9-\min(1.2,1.6d/L)$ for segment midpoints, where $L$ is the indexed duration (at least one second).
Global retrieval assigns $+1$ or $-1$ for overlap or non-overlap with an explicit time reference.
Absent applicable temporal cues, the temporal term is zero.
These fixed feature scales define the ranking heuristic; unit outer weights do not imply that the features contribute equally.

\subsection{Routing and Gate Feedback}
\label{sec:routing-details}

The first-tool selector applies ordered case-insensitive rules to the question.
Count, color, appearance, and on-screen-text cues take precedence and select $FI$.
Otherwise, causal or plot cues such as \texttt{why}, \texttt{cause}, and \texttt{throughout} select $GR$; temporal cues such as \texttt{when}, \texttt{before}, and \texttt{after} select $SR$; object or spatial cues such as \texttt{where} and \texttt{holding} select $FR$.
Unmatched questions default to $GR$.
This entry-point selector reads question text rather than dataset category labels; subsequent tool choices remain LLM-driven.

Separately, the implementation accepts optional question-type and time-reference context for verification policy hints.
When supplied, entity-recognition or key-information tags require targeted frame inspection, and explicit time references constrain admissible evidence.
This optional policy context is distinct from the first-tool selector.
A rejected answer appends a tool response and corrective guidance to the message history, so the next iteration receives an updated state rather than repeating an unchanged decision.

\section{Insights, Limitations, and Future Directions}
\label{sec:supp-limitations}

Finally, three challenging question types---dense counting, machine-capacity estimation, and transient color recognition---reveal a common principle: reliable long-video reasoning requires the granularity of the representation and verification process to match the granularity of the requested evidence.

\paragraph{Numerical reasoning as evidence aggregation.}
Counts and capacities are not always intrinsic attributes of a single visible entity.
They may require stable instance grouping across frames and a semantic distinction between visually salient instances and functionally valid units.
These requirements motivate count-aware visual queries, cross-frame agreement, and repeated narrow-window confirmation rather than a single unconstrained visual estimate.

\paragraph{State-aware graph abstraction.}
Entity and action relations provide strong localization cues, but some questions depend on latent numerical states such as occupancy, multiplicity, or capacity.
Enriching the hierarchy with temporally grounded state and cardinality relations would let graph retrieval propose not only where the relevant process occurs, but also which state must be verified.
Confidence on these relations can further determine when raw-frame confirmation remains necessary.

\paragraph{Question-adaptive temporal resolution.}
Fixed event and segment windows provide efficient long-range coverage, whereas brief attributes may change on a much shorter timescale.
Transient attributes therefore motivate uncertainty-driven temporal refinement: coarse retrieval first identifies the broader event, after which explicit time cues and state changes guide progressively narrower $FR$ or $FI$ queries.
These observations point to opportunities for extending, rather than replacing, the hierarchical design.
We continue to investigate instance-aware entity alignment, explicit count and occupancy relations, adaptive temporal splitting, and consensus-based visual verification.
Together, these directions aim to turn the hierarchy from a multi-scale retrieval index into an evidence-adaptive reasoning substrate whose resolution changes with the question.

